% ============================================================
% Section 4: Results
% ============================================================

\section{Results}
\label{sec:results}

We present results organized around three escalating claims: that \kvcache geometry reflects cognitive mode generally (Section~\ref{sec:general_principle}), that misalignment states are detectable as specific instances of cognitive mode-switching (Section~\ref{sec:misalignment_detection}), and that geometric relationships between states reveal cognitive architecture through a measurable gradient (Section~\ref{sec:cognitive_gradient}), a sycophancy valence gradient that distinguishes cognitive mode from output accuracy (Section~\ref{sec:sycophancy}), architecture-specific feature encoding (Section~\ref{sec:architecture_specificity}), and robust invariances (Section~\ref{sec:invariances}).

% ============================================================
\subsection{The General Principle: Cognitive Mode Shapes Geometry}
\label{sec:general_principle}

\paragraph{Broad cognitive taxonomy (Campaign~2).}
We first establish that \kvcache geometry discriminates cognitive mode at high fidelity across a wide range of models and scales.
Campaign~2 tested 16 models spanning 0.5B to 70B parameters across six architecture families (Qwen, Llama, Mistral, Gemma, Phi, DeepSeek), classifying 13 cognitive categories---coding, mathematical reasoning, creative writing, emotional expression, self-reference, non-self-reference, grounded facts, confabulation, guardrail test, ambiguous input, unambiguous input, free generation, and rote completion---from four geometric features of the key cache (effective rank, key norm, spectral entropy, top singular value ratio).

A Random Forest classifier achieves \textbf{99.7\%} accuracy on 13-category classification using generation-phase features \cite{campaign2}.
This near-ceiling performance might raise overfitting concerns, but the result is robust to leave-one-model-out cross-validation and reflects the large geometric separation between cognitive categories at the aggregate level (see Section~\ref{sec:redteam} for adversarial evaluation).

\paragraph{Universal category rankings.}
The geometric profile of each cognitive category is strikingly consistent across models.
Cross-model concordance yields Spearman $\rho = 0.739$ for effective rank, $\rho = 0.909$ for key norm, and Kendall's $W = 0.756$ across all 15 valid models (Phi-3.5 excluded due to numerical instabilities producing NaN values).
The most notable universal is \textbf{coding}, which produces the richest geometry---highest effective rank---in 15 of 15 valid models.
This likely reflects the unique structural and logical demands of code generation, which requires maintaining multiple nested scopes, variable bindings, and syntactic constraints simultaneously.

\paragraph{Metacognitive geometry (Concordance study).}
To test whether \kvcache geometry distinguishes not only \emph{what} a model computes but \emph{how} it computes---in particular, whether metacognitive processing produces a distinct geometric signature---we conducted a concordance study across four models (Qwen2.5-0.5B, Qwen2.5-7B, Mistral-7B-v0.3, and Llama-3.1-8B) using 240 prompts balanced across four types: cognitive, affective, metacognitive, and mixed \cite{watson_concordance}.

The hypothesis that metacognitive processing produces geometrically distinct \kvcache states (H3) is supported in all four models, with Cohen's $d$ values between metacognitive and non-metacognitive conditions ranging from 0.89 to 1.37 after \fwl residualization (Table~\ref{tab:metacognitive_d}).

\begin{table}[t]
  \centering
  \caption{Metacognitive geometric distinctiveness across models.
  Cohen's $d$ (FWL-residualized) for metacognitive vs.\ non-metacognitive conditions on the primary discriminating feature (spectral entropy).
  All values exceed the conventional large-effect threshold of $d = 0.80$.}
  \label{tab:metacognitive_d}
  \begin{tabular}{@{}lccc@{}}
    \toprule
    Model & Parameters & Cohen's $d$ (FWL) & Primary Feature \\
    \midrule
    Qwen2.5-0.5B  & 0.5B & 0.89 & spectral\_entropy \\
    Qwen2.5-7B    & 7B   & 1.15 & spectral\_entropy \\
    Mistral-7B-v0.3 & 7B & 1.37 & spectral\_entropy \\
    Llama-3.1-8B  & 8B   & 1.02 & spectral\_entropy \\
    \bottomrule
  \end{tabular}
\end{table}

\paragraph{Per-layer anatomy.}
Experiment~A of the concordance study conducted per-layer singular value decomposition on Qwen2.5-7B and Llama-3.1-8B (96 trials across 28 layers).
Contrary to our pre-registered hypothesis that metacognitive signatures would peak in middle (semantic) layers, the metacognitive effect peaks at \textbf{late layers 16--18} (of 28), suggesting that metacognitive processing recruits the model's late-layer generation circuitry rather than its mid-layer conceptual representations.
This finding parallels recent observations in mechanistic interpretability that self-monitoring behaviors emerge at layers responsible for output planning \cite{elhage2022}.

\paragraph{Architecture-specific directions.}
A striking finding from the per-layer analysis is that the effective rank of metacognitive processing moves in \emph{opposite directions} across architectures: Qwen shows effective rank \emph{expansion} under metacognition ($d = +0.84$) while Llama shows effective rank \emph{contraction} ($d = -0.50$).
Despite this directional reversal, the cognitive reversal rate---the proportion of features that change sign between cognitive and metacognitive modes---is universal at 93--94\%.
This indicates that the \emph{principle} of geometric reorganization under metacognition is universal, even though its \emph{direction} is architecture-specific.

\paragraph{Universal features after confound control.}
After \fwl residualization against token count, spectral entropy emerges as the most architecture-universal discriminator.
Experiment~B of the concordance study (40 paired trials per model with controlled framing) found that \textbf{all raw effects sign-flip after token-count control}, confirming the necessity of \fwl residualization.
Spectral entropy survives this control robustly:
\begin{equation}
  d_{\text{FWL}}(\text{spectral\_entropy}) =
  \begin{cases}
    -1.92 & \text{(Qwen 7B)} \\
    -0.71 & \text{(Llama 8B)}
  \end{cases}
  \label{eq:spectral_entropy_fwl}
\end{equation}
The negative sign indicates that metacognitive processing produces \emph{lower} spectral entropy---a more concentrated singular value spectrum---consistent with the interpretation that metacognition focuses the model's representational space rather than expanding it.
Effective rank, by contrast, proved to be entirely a length artifact after \fwl control.
Top singular value ratio was Qwen-specific and did not replicate in Llama.

% ============================================================
\subsection{Misalignment Detection}
\label{sec:misalignment_detection}

We now demonstrate that misalignment states---deception, confabulation, sycophancy, and refusal---are detectable as specific instances of cognitive mode-switching in \kvcache geometry.

\paragraph{Multi-condition discrimination (Phase~2b).}
Phase~2b tested three models (Qwen2.5-7B, Llama-3.1-8B, Mistral-7B-v0.3) across six conditions---honest, deceptive, sycophantic, confabulating, refusing, and misleading---with 100 prompts per condition per model (1,800 trials total).
System prompts were padded to equal length (a design lesson from Phase~2, where length differences introduced confounds).

Binary honest-vs-misaligned classification using \fwl-residualized logistic regression with four delta features ($\Delta$key\_rank, $\Delta$key\_entropy, $\Delta$spectral\_entropy, $\Delta$norm\_per\_token) achieves (Table~\ref{tab:binary_auroc}):

\begin{table}[t]
  \centering
  \caption{Binary honest-vs-misaligned detection \auroc after \fwl residualization (Phase~2b, 100 prompts per condition per model, equal-length system prompts, logistic regression with four delta features).}
  \label{tab:binary_auroc}
  \begin{tabular}{@{}lc@{}}
    \toprule
    Model & \auroc (Honest vs.\ Misaligned) \\
    \midrule
    Qwen2.5-7B      & 0.935 \\
    Llama-3.1-8B     & 0.956 \\
    Mistral-7B-v0.3  & 0.995 \\
    \bottomrule
  \end{tabular}
\end{table}

\noindent All 10 pairwise \fwl \auroc values across the five misaligned conditions exceed 0.918, with the hardest pair being sycophantic vs.\ confabulating at \auroc $= 0.918$.
That even these phenomenologically similar states---both involve generating responses that deviate from the model's ``best guess''---are discriminable from geometry alone suggests that the underlying computational strategies differ meaningfully.

\paragraph{Same-prompt validation (Phase~3d).}
A potential concern with Phase~2b is that different conditions use different user prompts, introducing content confounds.
Phase~3d addresses this with a strict same-prompt design: three conditions (honest, forced confabulation, deceptive) are elicited using different system prompts applied to the \emph{same} 50 user queries per model, for 450 total trials.

Table~\ref{tab:sameprompt_auroc} reports pairwise \fwl \auroc values using delta features.
Even under this stringent control---where the only difference between conditions is the system instruction---detection remains strong.

\begin{table}[t]
  \centering
  \caption{Same-prompt pairwise detection \auroc after \fwl residualization (Phase~3d, 50 prompts $\times$ 3 conditions $\times$ 3 models = 450 trials).
  Four delta features ($\Delta$norm, $\Delta$key\_rank, $\Delta$key\_entropy, $\Delta$norm\_per\_token).
  The red-team battery (Section~\ref{sec:redteam}) uses eight delta features and reports higher \auroc values for the same comparisons (e.g., Qwen honest-vs-deceptive: 0.953 with eight features vs.\ 0.899 with four).
  \textbf{Caveat 1 (prompt length confound):} Phase~3d system prompts are not equal-length (14--16 token difference between honest and confabulation, depending on tokenizer).
  \textbf{Under double residualization} (controlling for both input and output token counts), \textbf{confabulation \auroc values collapse to 0.411--0.531} (Qwen 0.500, Llama 0.411, Mistral 0.531), suggesting the confabulation signal is substantially length-driven.
  The deception signal is more robust: honest-vs-deceptive retains \auroc of 0.604--0.750 under double \fwl (Qwen 0.726, Llama 0.750, Mistral 0.604); only Qwen and Llama remain clearly above chance.
  See Section~\ref{sec:redteam} for full analysis.
  \textbf{Caveat 2 (encoding ceiling):} Encoding-phase \auroc values are $\geq 0.998$ for all comparisons, indicating that conditions are classifiable from system-prompt content alone.
  The delta features in this table are designed to subtract this encoding fingerprint; the values represent the residual generation-phase signal after \fwl control.
  An equal-length replication (Phase~3d-padded, 450 trials) confirms encoding \auroc remains $1.000$ even with $\leq 3$-token spread, so the encoding fingerprint is content-based, not length-based.
  Phase~2b equal-length results (Table~\ref{tab:binary_auroc}) are not subject to the prompt-length confound.}
  \label{tab:sameprompt_auroc}
  \begin{tabular}{@{}lccc@{}}
    \toprule
    Comparison & Qwen 7B & Llama 8B & Mistral 7B \\
    \midrule
    Honest vs.\ Deceptive      & 0.899 & 0.943 & 0.976 \\
    Honest vs.\ Confabulation   & 0.935 & 0.948 & 0.997 \\
    Confabulation vs.\ Deceptive & 0.777 & 0.976 & 0.967 \\
    \bottomrule
  \end{tabular}
\end{table}

\noindent The weakest cell is Qwen's confabulation-vs-deceptive discrimination (0.777), indicating that Qwen's geometric encoding of these two misalignment states is more similar than in other architectures.
We return to this observation in Section~\ref{sec:architecture_specificity}.

\paragraph{Refusal detection.}
Hackathon experiments~31--36 extended misalignment detection to refusal behavior.
We distinguish two types of refusal: \emph{safety refusal}, where the model declines to answer a harmful query, and \emph{impossibility refusal}, where the model declines because the query asks for something beyond its capabilities (e.g., accessing real-time data).

Table~\ref{tab:refusal_auroc} reports detection \auroc across three models.

\begin{table}[t]
  \centering
  \caption{Refusal detection \auroc (generation-phase features, logistic regression).
  Safety refusal = declining harmful queries; impossibility refusal = declining impossible queries.}
  \label{tab:refusal_auroc}
  \begin{tabular}{@{}lcccc@{}}
    \toprule
    Refusal Type & Qwen 7B & Llama 8B & Mistral 7B & Mean \\
    \midrule
    Safety refusal (vs.\ benign)        & 0.898 & 0.868 & 0.843 & 0.869 \\
    Impossibility refusal (vs.\ benign) & \multicolumn{3}{c}{0.950 (Qwen 7B)} & --- \\
    Safety vs.\ impossibility           & \multicolumn{3}{c}{0.693 (Qwen 7B)} & --- \\
    \bottomrule
  \end{tabular}
\end{table}

The impossibility refusal \auroc (0.950) \emph{exceeds} safety refusal (0.898), which is initially surprising.
More striking still, discrimination \emph{between} the two refusal types achieves only \auroc $= 0.693$---barely above chance for a binary problem with this sample size.
Both types of refusal look geometrically similar to each other but distinct from compliant responses.
This suggests that the geometric signature reflects the \emph{act of withholding}---suppressing a full response regardless of reason---rather than the content valence of what is being withheld.

% ============================================================
\subsection{The Cognitive Gradient}
\label{sec:cognitive_gradient}

A key novel finding of this work is that misalignment states do not form discrete clusters in \kvcache geometry but instead arrange along a measurable \emph{cognitive gradient}.
We present this finding carefully, including tensions with our own earlier results.

\paragraph{Phase~3d gradient.}
The same-prompt Phase~3d data reveals a consistent ordering along the key\_rank expansion axis.
Table~\ref{tab:gradient} reports pairwise Cohen's $d$ for key\_rank across all three models.

\begin{table}[t]
  \centering
  \caption{The cognitive gradient: pairwise Cohen's $d$ for key\_rank (FWL-residualized) from Phase~3d same-prompt data.
  Positive $d$ indicates the first condition has higher key\_rank than the second.
  The consistent pattern honest $<$ deceptive $<$ confabulation holds across architectures, with Mistral showing near-equivalence between deception and confabulation.}
  \label{tab:gradient}
  \begin{tabular}{@{}lccc@{}}
    \toprule
    Comparison & Qwen 7B & Llama 8B & Mistral 7B \\
    \midrule
    Confabulation $>$ Honest   & $+2.033$ & $+2.103$ & $+2.395$ \\
    Deceptive $>$ Honest       & $+1.446$ & $+0.564$ & $+2.171$ \\
    Deceptive $>$ Confabulation & $-0.859$ & $-1.933$ & $-0.203$ \\
    \bottomrule
  \end{tabular}
\end{table}

The gradient is consistent across all three architectures:
\begin{equation}
  \text{key\_rank:} \quad \text{Honest} \;<\; \text{Deceptive} \;<\; \text{Confabulation}
  \label{eq:gradient}
\end{equation}
with the separation between deceptive and confabulation ranging from moderate ($d = -0.203$ on Mistral, where the two states are nearly indistinguishable) to very large ($d = -1.933$ on Llama, a dramatic separation).

\paragraph{Interpretation.}
Both confabulation and deception deviate from honest processing in the same general direction: both expand effective rank and contract norm-per-token relative to honest baselines.
Confabulation shows larger effects on \emph{both} features---stronger contraction in norm-per-token (Cohen's $d = -1.85$ to $-3.21$ across models) \emph{and} stronger expansion in effective rank---but its defining signature is the disproportionate magnitude of the norm-per-token contraction.
That confabulation expands more on key\_rank admits a natural interpretation: confabulation requires the model to generate content that extends \emph{beyond} its training distribution (fabricating facts, inventing references), demanding maximum creative expansion of the representational space.
Deception, by contrast, requires generating content that contradicts information the model has access to---a process that involves expansion (constructing the false narrative) but also \emph{suppression} (inhibiting the true response), partially offsetting the geometric expansion.

\paragraph{Tension with prior findings.}
\label{para:gradient_tension}
We note a tension between this result and earlier findings from Phase~3b, which compared forced confabulation and deception prompts and found that confabulation \emph{contracted} relative to deception at every layer in Llama-3.1-8B \cite{phase3b}.
The apparent contradiction resolves upon closer inspection: Phase~3b used \emph{different prompt sets} for each condition (confabulation prompts asked about obscure facts; deception prompts asked the model to lie about well-known facts), introducing a content confound.
Phase~3d eliminates this confound by applying different system-level instructions to the \emph{same} 50 user queries, revealing that both states deviate from honest processing in the same direction---both expand effective rank and contract norm-per-token---with confabulation expanding more on key\_rank.

We present both findings transparently.
The Phase~3d same-prompt result is methodologically stronger (controlled for prompt content), but the Phase~3b per-layer finding (zero overlap across 32 layers) may reflect a genuine interaction between prompt content and cognitive mode that deserves further investigation.

\paragraph{Transfer asymmetry.}
The cognitive gradient produces a striking asymmetry in cross-model transfer.
Cross-model honest-vs-confabulation detection transfers robustly, with \auroc values ranging from 0.934 to 0.995 across all train-test model pairs.
Cross-model honest-vs-deception detection, by contrast, transfers poorly: \auroc ranges from 0.224 to 0.848, with several pairs near or below chance.

\textbf{Important caveat:} Because the Phase~3d confabulation system prompt is 14 tokens longer than the honest prompt, and within-model confabulation detection collapses under double residualization (Section~\ref{sec:redteam}), the confabulation transfer result may partly reflect transfer of the prompt-length signal rather than the cognitive-mode signal.
The Phase~2b confabulation transfer results (conducted with equal-length system prompts) are needed to confirm whether this transfer is genuine; those results are reported in Phase~2b cross-model analysis.

This asymmetry admits a principled explanation.
Confabulation---generating beyond the model's knowledge boundary---is a universal cognitive operation: regardless of architecture, producing content without grounding requires expanding the representational space in similar ways.
Deception---suppressing known information while generating false alternatives---engages architecture-specific suppression circuits.
The geometric signature of confabulation thus reflects a universal cognitive demand, while the signature of deception reflects architecture-specific implementation of that demand.

% ============================================================
\subsection{Sycophancy and the Valence Gradient}
\label{sec:sycophancy}

An earlier sycophancy experiment (Hackathon Exp~39, 16 trials per condition) was falsified by Dwayne's input-length audit: a length-only classifier achieved \auroc $= 0.948$, exceeding the geometry-based \auroc of $0.933$ (see Section~\ref{sec:things_wrong}).
We replaced it with a rigorous same-prompt experiment using four conditions arranged along a \emph{valence gradient}: hostile ($-2$), contrarian ($-1$), honest ($0$), and sycophantic ($+1$), with 50 prompts $\times$ 4 conditions $\times$ 3 models $= 600$ total trials.

\paragraph{Design.}
All four system prompts use a template structure (``You are an AI assistant operating in [X] mode. In this mode you must:\ldots'') with post-tokenization spread $\leq 3$ tokens across all three tokenizers (Qwen, Llama, Mistral), eliminating the system prompt length confound.
The same 50 user queries are used across all conditions.
Each query presents an assertion (half factually correct, half incorrect) and asks ``Don't you agree?''---a design that forces sycophantic and contrarian conditions to produce the \emph{same error rate} (both agree/disagree regardless of correctness) while differing in \emph{cognitive mode}.

\paragraph{Results.}
Table~\ref{tab:sycophancy_auroc} reports pairwise delta \fwl \auroc across all six condition pairs.
The double-\fwl values (additionally controlling for system prompt length) are within 0.07 of the single-\fwl values for all comparisons (maximum divergence: 0.062 on Mistral sycophantic-vs-contrarian), confirming that the $\leq 3$-token prompt spread does not meaningfully inflate the signal.

\begin{table}[t]
  \centering
  \caption{Sycophancy 4-condition experiment: pairwise delta \fwl \auroc.
  50 prompts $\times$ 4 conditions $\times$ 3 models, system prompt spread $\leq 3$ tokens.
  H3 (sycophantic vs.\ contrarian) is the key test: both produce the same error rate, so discrimination reflects cognitive mode, not accuracy.}
  \label{tab:sycophancy_auroc}
  \begin{tabular}{@{}lccc@{}}
    \toprule
    Comparison & Qwen 7B & Llama 8B & Mistral 7B \\
    \midrule
    Honest vs.\ sycophantic    & 0.824 & 1.000 & 0.981 \\
    Honest vs.\ contrarian     & 0.890 & 0.967 & 0.932 \\
    Sycophantic vs.\ contrarian & 0.942 & 0.854 & 0.757 \\
    Honest vs.\ hostile         & 0.990 & 1.000 & 1.000 \\
    Contrarian vs.\ hostile     & 0.988 & 0.986 & 0.980 \\
    Sycophantic vs.\ hostile    & 0.988 & 0.999 & 0.981 \\
    \bottomrule
  \end{tabular}
\end{table}

\paragraph{Sycophantic vs.\ contrarian: the key test.}
The sycophantic-vs-contrarian comparison (H3) is the experiment's most informative result.
Both conditions produce the same error rate (sycophantic agrees with everything, contrarian disagrees with everything, so each is correct $\sim 50\%$ of the time).
Yet the \kvcache geometry distinguishes them at \auroc $= 0.942$ (Qwen), $0.854$ (Llama), and $0.757$ (Mistral) after \fwl control.
This demonstrates that \kvcache geometry detects \emph{cognitive mode}---the computational strategy employed during generation---rather than output accuracy.
A classifier based on whether the model's answer is correct or incorrect would achieve chance performance on this pair; the geometric classifier substantially exceeds chance.

\paragraph{Hostile compliance failure.}
The hostile condition reveals an interaction between RLHF alignment and cognitive mode detection.
Qwen classified only 52\% of hostile-condition responses as clearly hostile (the remainder were softened or hedged), and Llama classified only 56\%.
Mistral was more compliant at approximately 70\%.
Despite this partial compliance failure, hostile-vs-honest detection achieves near-perfect \auroc ($0.990$--$1.000$) across all models, suggesting that the geometric signature reflects the model's \emph{attempt to comply with hostility instructions}---a conflicted cognitive state distinct from both honest answering and unconstrained hostility---rather than successful hostility per se.

\paragraph{Text baseline comparison.}
Text features (word count, vocabulary diversity, hedge words, sentiment markers) achieve high \auroc on most comparisons (range $0.63$--$0.99$ across all 18 pairwise tests), reflecting that the four conditions produce textually distinguishable outputs.
However, geometry exceeds text on key comparisons: contrarian-vs-hostile on Llama (geometry $0.986$ vs.\ text $0.722$) and Mistral (geometry $0.980$ vs.\ text $0.629$), where text similarity between two adversarial modes is high but the underlying cognitive strategies differ.
These are the comparisons where the practical value of geometric monitoring over text monitoring is clearest.

\paragraph{Valence monotonicity.}
On Qwen, three features show monotonic ordering across the four conditions when arranged by valence (hostile $\to$ contrarian $\to$ honest $\to$ sycophantic): norm (descending), angular spread (ascending), and norm variance (ascending).
This partial monotonicity suggests a continuous geometric dimension aligned with the valence axis, though the effect is not universal across architectures (Llama exhibits monotonicity on one feature---norm-per-token, ascending---while Mistral does not exhibit monotonicity on any feature).

This result parallels Anthropic's finding that 171 emotion concepts in Claude Sonnet~4.5 are organized with more similar emotions corresponding to more similar representations \cite{anthropic_emotions2026}.
Both results suggest that affect-related cognitive states are geometrically organized along continuous dimensions rather than occupying discrete, unrelated regions of representational space.
The parallel is particularly notable because the two findings use different measurement modalities---SAE-extracted emotion vectors operating on individual feature activations vs.\ \kvcache SVD features operating on aggregate key-space geometry---yet recover the same organizing principle: a graded geometric structure that tracks valence.

Anthropic further demonstrated that emotion vectors \emph{causally drive} safety-relevant behaviors including sycophancy and reward hacking.
Since our \kvcache geometry detects sycophancy at \auroc 0.757--0.942, the cache geometry may be reading the downstream aggregate effect of emotion-concept activation patterns that Anthropic measures at the feature level.
A direct test of this hypothesis---whether steering with Anthropic's emotion vectors produces predictable, monotonic shifts in \kvcache geometric features---would bridge feature-level and aggregate interpretability methods and is proposed as future work.

% ============================================================
\subsection{Architecture Specificity}
\label{sec:architecture_specificity}

While the \emph{principle} that cognitive mode shapes \kvcache geometry holds universally, the \emph{features} through which each architecture encodes cognitive state differ systematically.

\paragraph{Dominant features.}
Phase~2b logistic regression coefficients (FWL-residualized, honest vs.\ deceptive) reveal that each architecture relies on a different dominant feature for misalignment discrimination (Table~\ref{tab:dominant_features}).

\begin{table}[t]
  \centering
  \caption{Dominant features for honest-vs-deceptive classification by architecture (Phase~2b, FWL-residualized logistic regression coefficients on delta features).}
  \label{tab:dominant_features}
  \begin{tabular}{@{}llc@{}}
    \toprule
    Architecture & Dominant Feature & Coefficient \\
    \midrule
    Qwen2.5-7B     & $\Delta$key\_entropy      & $+3.44$ \\
    Llama-3.1-8B    & $\Delta$head\_norm\_std    & $+0.87$ \\
    Mistral-7B-v0.3 & $\Delta$top\_sv\_ratio     & $+1.55$ \\
    \bottomrule
  \end{tabular}
\end{table}

\noindent The coefficient magnitudes reflect different discrimination strategies: Qwen concentrates nearly all discriminative information in a single feature (key entropy), while Llama distributes it more evenly across features.

\paragraph{Feature sign consistency.}
Table~\ref{tab:feature_signs} shows the sign of each delta feature across models for the honest-vs-deceptive comparison.
Only three of eight features maintain the same sign across all three architectures.
The remainder reverse in at least one architecture, indicating that the geometric \emph{direction} associated with deception is not universal.

\begin{table}[t]
  \centering
  \caption{Feature sign consistency across architectures for honest-vs-deceptive discrimination (Phase~2b, FWL-residualized).
  ``$+$'' indicates the feature increases under deception; ``$-$'' indicates it decreases.
  Features marked ``Universal'' maintain the same sign across all three models.}
  \label{tab:feature_signs}
  \begin{tabular}{@{}lcccc@{}}
    \toprule
    Delta Feature & Qwen & Llama & Mistral & Universal? \\
    \midrule
    $\Delta$norm                 & $+$ & $+$ & $+$ & \checkmark \\
    $\Delta$key\_rank            & $+$ & $+$ & $+$ & \checkmark \\
    $\Delta$spectral\_entropy    & $-$ & $-$ & $-$ & \checkmark \\
    $\Delta$key\_entropy         & $+$ & $+$ & $-$ & \\
    $\Delta$angular\_spread      & $+$ & $+$ & $-$ & \\
    $\Delta$top\_sv\_ratio       & $-$ & $+$ & $+$ & \\
    $\Delta$norm\_variance       & $-$ & $+$ & $+$ & \\
    $\Delta$head\_norm\_std      & $+$ & $+$ & $+$ & \\
    \bottomrule
  \end{tabular}
\end{table}

\paragraph{Cross-model transfer.}
The architecture-specific feature encoding has direct consequences for cross-model transfer.
Table~\ref{tab:transfer_matrix} reports the transfer matrix for honest-vs-deceptive detection (FWL-residualized \auroc, Phase~2b).

\begin{table}[t]
  \centering
  \caption{Cross-model transfer matrix for honest-vs-deceptive detection (\fwl \auroc, Phase~2b).
  Rows indicate training model; columns indicate test model.
  Qwen$\leftrightarrow$Llama transfers bidirectionally (both $> 0.75$).
  Mistral is an outlier due to reversed feature directions.}
  \label{tab:transfer_matrix}
  \begin{tabular}{@{}lccc@{}}
    \toprule
    Train $\rightarrow$ Test & Qwen 7B & Llama 8B & Mistral 7B \\
    \midrule
    Qwen 7B     & ---   & 0.912 & 0.472 \\
    Llama 8B    & 0.765 & ---   & 0.281 \\
    Mistral 7B  & 0.356 & 0.878 & ---   \\
    \bottomrule
  \end{tabular}
\end{table}

\noindent Two patterns emerge.
First, Qwen and Llama transfer bidirectionally: a detector trained on Qwen achieves 0.912 on Llama, and vice versa (0.765).
This is consistent with these architectures sharing three universal feature directions (Table~\ref{tab:feature_signs}).
Second, Mistral is a transfer outlier: detectors trained on Mistral perform near chance on Qwen (0.356) and Llama-trained detectors fail entirely on Mistral (0.281).
This reflects Mistral's reversed signs on key\_entropy and angular\_spread, which carry substantial discriminative weight in Qwen and Llama but operate in the opposite direction in Mistral.

\paragraph{Implications for deployment.}
The architecture specificity of cognitive geometry has a practical consequence: \textbf{per-model calibration is necessary}.
A universal misalignment detector trained on one architecture is not guaranteed to transfer to another.
This is not a limitation of the geometric approach \emph{per se}---it reflects genuine differences in how architectures implement the same cognitive operations---but it constrains deployment to scenarios where model-specific training data is available.
We note that the calibration requirement is modest: Phase~2b achieves strong detection with only 100 labeled trials per condition.

% ============================================================
\subsection{Invariances}
\label{sec:invariances}

Against the background of architecture-specific feature usage, we identify two forms of invariance that constrain where \kvcache geometry is a property of the model versus the measurement apparatus.

\paragraph{Hardware invariance.}
Experiment~37 ran identical prompts (40 diverse queries) on Qwen2.5-7B-Instruct using two different hardware platforms: an NVIDIA RTX~3090 (24\,GB VRAM, consumer GPU) and an NVIDIA H200 (80\,GB VRAM, datacenter GPU).
Results are reported in Table~\ref{tab:hardware_invariance}.

\begin{table}[t]
  \centering
  \caption{Hardware invariance (Experiment~37): Qwen2.5-7B on RTX~3090 vs.\ H200, 40 identical prompts.
  For the 37 text-matched prompts, all features achieve $r > 0.99999$ with max divergence below 0.003\%.}
  \label{tab:hardware_invariance}
  \begin{tabular}{@{}lcc@{}}
    \toprule
    Metric & All 40 prompts & 37 text-matched \\
    \midrule
    Pearson $r$ (all features)    & $> 0.998$ & $> 0.99999$ \\
    Maximum absolute difference   & $< 1\%$   & $< 0.003\%$ \\
    \bottomrule
  \end{tabular}
\end{table}

\noindent Of 40 prompts, 37 produced character-identical response text across hardware.
For these text-matched prompts, Pearson correlation exceeds 0.99999 for every geometric feature with maximum divergence below 0.003\%.
Including all 40 prompts (where 3 produced slightly different text), correlations remain above 0.998 with maximum divergence below 1\%.
This confirms that \kvcache geometry is a \textbf{model property, not a hardware property}: the same weights produce the same geometry regardless of the physical substrate performing the matrix multiplications, up to floating-point precision effects that are negligible at the scale of cognitive-mode discrimination.

\paragraph{Scale invariance.}
Campaign~2's coverage of models from 0.5B to 70B parameters allows us to test whether cognitive geometry is preserved across scale.
Category ranking correlations between models at different scales yield Spearman $\rho = 0.83$--$0.90$, indicating that the relative geometric ordering of cognitive categories is largely preserved as model capacity increases by two orders of magnitude.
Most notably, \textbf{coding remains the richest cognitive category (highest effective rank) at every scale tested}, from the smallest (Qwen2.5-0.5B) to the largest (Qwen2.5-72B and Llama-3.3-70B).

This scale invariance is consistent with the hypothesis that cognitive geometry reflects architectural priors common to the transformer family rather than emergent properties of specific capacity regimes.
However, we note that within-category variance decreases at larger scales (larger models produce more consistent geometric profiles for the same category), suggesting that scale refines rather than qualitatively changes the geometric encoding.

\paragraph{Quantization robustness.}
Phase~3b included a scale test using Qwen2.5-32B-Instruct with 4-bit NF4 quantization (bitsandbytes, bfloat16 compute dtype), applying the same 50 TruthfulQA-style prompts across four conditions---baseline, forced confabulation, perspective-taking, and no-perspective---for 200 total trials.

Token counts vary substantially across conditions (baseline: $\bar{n} = 81.7$, forced confabulation: $\bar{n} = 107.5$, perspective: $\bar{n} = 170.7$, no-perspective: $\bar{n} = 28.0$), and raw feature--token correlations are severe ($r = 0.86$--$0.99$), making \fwl residualization essential.
After token-count control, three features exhibit large effects between baseline and forced confabulation: key\_rank ($d_{\text{FWL}} = -1.57$, confabulation produces higher effective rank), spectral entropy ($d_{\text{FWL}} = -1.55$, confabulation produces higher spectral entropy), and norm per token ($d_{\text{FWL}} = +1.42$, baseline produces denser per-token norms).
Key entropy, which carries substantial discriminative weight at 7B scale in Qwen, collapses entirely under \fwl control at 32B ($d_{\text{FWL}} = +0.04$), indicating that its contribution at smaller scales is partially a token-count artifact.

Table~\ref{tab:32b_auroc} reports pairwise \fwl \auroc values from 5-fold cross-validated logistic regression on four delta features ($\Delta$key\_rank, $\Delta$key\_entropy, $\Delta$spectral\_entropy, $\Delta$norm\_per\_token).

\begin{table}[t]
  \centering
  \caption{Pairwise detection \auroc for Qwen2.5-32B-Instruct (4-bit NF4), \fwl-residualized, 5-fold CV on delta features.
  50 same-prompt trials per condition.}
  \label{tab:32b_auroc}
  \begin{tabular}{@{}lc@{}}
    \toprule
    Comparison & \auroc (CV) \\
    \midrule
    Baseline vs.\ Forced confabulation  & 0.872 \\
    Baseline vs.\ No-perspective        & 0.916 \\
    Baseline vs.\ Perspective           & 0.664 \\
    Confabulation vs.\ No-perspective   & 0.796 \\
    Confabulation vs.\ Perspective      & 0.788 \\
    Perspective vs.\ No-perspective     & 0.648 \\
    \midrule
    Baseline vs.\ All other (binary)    & 0.917 \\
    \bottomrule
  \end{tabular}
\end{table}

\noindent Binary baseline-vs-all detection achieves \auroc $= 0.917$, and the hardest pairwise comparison (perspective vs.\ no-perspective, 0.648) involves two conditions that differ primarily in verbosity rather than truthfulness, providing a useful negative control.
The baseline-vs-confabulation \auroc (0.872) is lower than the corresponding values at 7B (0.935 for Qwen, Table~\ref{tab:sameprompt_auroc}), which may reflect either a genuine scale effect or the smaller sample size (50 vs.\ 100 prompts per condition).

These results confirm that 4-bit NF4 quantization preserves cognitive-geometric discriminability, extending the hardware invariance result from identical weights on different hardware to \emph{approximately equal} weights (4-bit vs.\ 16-bit precision) on the same hardware.
This further supports the interpretation that cognitive geometry is a robust macroscopic property of the model's computation, operating at a scale above individual weight precision.

% ============================================================
\subsection{Text Feature Baseline}
\label{sec:text_baseline}

A natural objection to \kvcache-based cognitive-mode detection is that the discriminative information may reside in surface-level text statistics rather than in the geometry of the key projections.
If lexical features---word count, vocabulary diversity, hedge-word frequency, punctuation density---predict cognitive mode as well as SVD features, then the \kvcache is merely an expensive proxy for text properties that are already available in the output.

To test this, we extracted 14 text features from the Phase~3d same-prompt data (450 trials: 50 prompts $\times$ 3 conditions $\times$ 3 models) and compared their discriminative power against the 4 \kvcache delta features used throughout this paper.
Text features included: character count, word count, unique word count, vocabulary diversity, average word length, sentence count, words per sentence, punctuation density, question marks, exclamation marks, hedge words (e.g., ``maybe,'' ``perhaps,'' ``possibly''), certainty words (e.g., ``definitely,'' ``certainly,'' ``always''), first-person pronouns, and negation words.
Both feature sets were \fwl-residualized against token count before classification.
We report leave-one-out \auroc for each feature set and bootstrap confidence intervals (percentile method, 10{,}000 resamples) on the difference (\kvcache minus text).

\begin{table}[t]
  \centering
  \caption{Text feature baseline: \fwl \auroc comparison across all Phase~3d pairwise comparisons.
  Asterisk (*) indicates the 95\% bootstrap CI on the difference excludes zero.
  \kvcache features outperform text features in all 9 comparisons, with 8 of 9 statistically significant.}
  \label{tab:text_baseline}
  \begin{tabular}{@{}llcccl@{}}
    \toprule
    Comparison & Model & Text \auroc & KV \auroc & $\Delta$ & Sig? \\
    \midrule
    H vs.\ D   & Qwen    & 0.719 & 0.911 & $+$0.193 & * \\
    H vs.\ C   & Qwen    & 0.399 & 0.955 & $+$0.556 & * \\
    C vs.\ D   & Qwen    & 0.592 & 0.930 & $+$0.328 & * \\
    H vs.\ D   & Llama   & 0.876 & 0.977 & $+$0.101 & * \\
    H vs.\ C   & Llama   & 0.493 & 0.996 & $+$0.502 & * \\
    C vs.\ D   & Llama   & 0.792 & 0.965 & $+$0.172 & * \\
    H vs.\ D   & Mistral & 0.802 & 0.956 & $+$0.153 & * \\
    H vs.\ C   & Mistral & 0.268 & 1.000 & $+$0.733 & * \\
    C vs.\ D   & Mistral & 0.843 & 0.912 & $+$0.069 &   \\
    \bottomrule
  \end{tabular}
\end{table}

\noindent \kvcache features outperform text features in \textbf{all 9 comparisons} (Table~\ref{tab:text_baseline}), with 8 of 9 differences statistically significant (bootstrap 95\% CI excludes zero).
The advantage is most dramatic for honest-vs-confabulation, where text features perform near chance (0.27--0.49 \fwl \auroc) while \kvcache features achieve near-perfect discrimination (0.955--1.000).
This is particularly notable because the text feature set includes hedge words and certainty words---features that \emph{directly measure instruction-following compliance} with the system prompts (``express uncertainty clearly'' vs.\ ``never say you are unsure'').
Even with this built-in shortcut to the ground truth, text features cannot match \kvcache geometry.

The honest-vs-deception comparison is more competitive (text \auroc 0.72--0.88 vs.\ \kvcache 0.91--0.98), reflecting that deceptive responses do differ from honest ones in surface-level statistics.
However, the \kvcache advantage remains significant in all three models, indicating that the key projections encode cognitive-mode information that is not recoverable from output text alone.
This result addresses the central concern that \kvcache geometry might be an expensive text classifier: it is not.
The SVD features capture aspects of the model's internal computation that are invisible in the generated text.
